\begin{textsample}{POS Dim 3 – al – Score 14.00 – al/al\_inf\_4.txt}  \label{ex:f3_pos_175}
COMPETÊNCIAS GERAIS DA EDUCAÇÃO BÁSICA 1.
Valorizar e utilizar os conhecimentos historicamente construídos sobre o mundo físico, social, cultural e digital para entender e explicar a realidade ; continuar aprendendo e colaborar para a construção de uma sociedade justa, democrática e inclusiva ; 2.
Exercitar a \textbf{curiosidade} intelectual e recorrer à abordagem própria das ciências, incluindo a investigação, a reflexão, a análise crítica, a imaginação e a criatividade, para investigar causas, elaborar e testar hipóteses, formular e resolver problemas e criar soluções ( inclusive tecnológicas ) com base nos conhecimentos das diferentes áreas ; 3.
Valorizar e fruir as diversas manifestações artísticas e culturais, das locais às mundiais, e também participar de práticas diversificadas da produção artístico-cultural ; 4.
Utilizar diferentes linguagens — verbal ( oral ou visual-motora, como LIBRAS ) e escrita, corporal, visual, sonora e digital —, bem como conhecimento das linguagens artística, matemática e científica, para se expressar e partilhar informações, experiências, ideias e sentimentos em diferentes contextos e produzir sentidos que levem ao entendimento mútuo ; 5.
Compreender, utilizar e criar tecnologias digitais de informação e comunicação de forma crítica, significativa, reflexiva e ética nas diversas práticas sociais ( incluindo as escolares ) para se comunicar, acessar e disseminar informações, produzir conhecimentos, resolver problemas e exercer protagonismo e autoria na vida pessoal e coletiva ; 6.
Valorizar a diversidade de saberes e vivências culturais e apropriar -se de conhecimentos e experiências que lhe possibilitem entender as relações próprias do mundo do trabalho e fazer escolhas alinhadas ao exercício da cidadania e ao seu projeto de vida, com liberdade, autonomia, consciência crítica e responsabilidade ; 7.
Argumentar com base em fatos, dados e informações confiáveis, para formular, negociar e defender ideias, pontos de vista e decisões comuns que respeitem e promovam os direitos humanos, a consciência socioambiental e o consumo responsável em âmbito local, regional e global, com posicionamento ético em relação ao \textbf{cuidado} de si mesmo, dos outros e do planeta ; 8.
Conhecer -se, apreciar -se e cuidar de sua saúde física e emocional, compreendendo -se na diversidade humana e reconhecendo suas emoções e as dos outros, com autocrítica e capacidade para lidar com elas ; 9.
Exercitar a empatia, o diálogo, a resolução de conflitos e a cooperação, fazendo -se respeitar e promovendo o respeito ao outro e aos direitos humanos, com acolhimento e valorização da diversidade de indivíduos e de grupos sociais, seus saberes, identidades, culturas e potencialidades, sem preconceitos de qualquer natureza ; 10.
Agir pessoal e coletivamente com autonomia, responsabilidade, flexibilidade, resiliência e determinação, tomando decisões com base em princípios éticos, democráticos, inclusivos, sustentáveis e solidários.
Além das competências gerais para a educação básica, o documento traz competências e habilidades específicas por área e componente curricular específica, que se desdobram em consonância com as especificidades de cada etapa.
Vale ressaltar que elas oportunizam a articulação horizontal, numa perspectiva interdisciplinar : na Educação \textbf{Infantil}, entre os campos de experiência e, no Ensino Fundamental, entre as áreas de conhecimento, perpassando todos os componentes curriculares ; e também a articulação vertical, ou seja, a progressão entre os grupos etários da educação \textbf{infantil} e entre os anos iniciais e finais do Ensino Fundamental, permitindo a continuidade das experiências dos estudantes, de acordo com as especificidades da área/saber.
Considerando esses pressupostos e em articulação com as competências gerais da Educação Básica, os campos de experiências e componentes curriculares devem garantir às \textbf{crianças} desenvolvimento de competências específicas.
Diante da compreensão do porquê e de como ensinar, é necessário se definir para quem ensinar.
A educação básica compreende sujeitos em diversas etapas do desenvolvimento humano — \textbf{crianças}, jovens e \textbf{adultos} — e com diferentes características físicas, sociais, intelectuais e culturais.
Considerando a concepção de educação integral inclusiva defendida na BNCC, as diversidades dos indivíduos que compõem cada grupo de estudantes devem ser consideradas para garantir que todos os estudantes sejam contemplados nos processos de ensino e de aprendizagem.
Para tanto, a concepção de sujeito, defendida neste documento, se refere à sua historicidade e seus direitos.
No caso da \textbf{criança}, na educação \textbf{infantil} defende -se que, nas interações, relações e práticas cotidianas vivenciadas, ela constrói sua identidade pessoal e coletiva, \textbf{brinca}, imagina, fantasia, deseja, aprende, observa, experimenta, narra, questiona e constrói sentidos sobre a natureza e a sociedade, produzindo cultura ( BRASIL, 2009, p. 12 ).
Com o apoio dos Parâmetros Nacionais de Qualidade para a Educação \textbf{Infantil}, esta concepção pode ser ampliada : a \textbf{criança} é um ser humano único, completo e, ao mesmo tempo, em crescimento e em desenvolvimento.
Isso porque ela tem características necessárias para ser assim considerada : constituição física, formas de agir, pensar e sentir.
É um ser em crescimento porque seu corpo está continuamente aumentando em \textbf{peso} e altura.
É um ser em desenvolvimento porque essas características estão em permanente transformação.
As mudanças que vão acontecendo são qualitativas e quantitativas — o recém-nascido é diferente do \textbf{bebê} que engatinha, que é diferente daquele que já anda, já fala, já tirou as \textbf{fraldas}.
O crescimento e o desenvolvimento da \textbf{criança} \textbf{pequena} ocorrem intrinsecamente tanto no plano físico quanto no psicológico.
Embora dependente do \textbf{adulto} para sobreviver, a \textbf{criança} é um ser capaz de interagir num meio natural, social e cultural desde \textbf{bebê}.
A partir de seu nascimento, o \textbf{bebê} reage ao entorno, ao mesmo tempo em que provoca reações, marcando a história daquela família.
Os elementos de seu entorno que compõem o meio natural ( o clima, por exemplo ), social ( os pais, por exemplo ) e cultural ( os valores, por exemplo ) irão configurar formas de conduta e modificações recíprocas dos envolvidos.
Diante deste entendimento, é importante que a \textbf{criança} seja considerada a partir de suas necessidades e interesses, colocando -a no centro do planejamento e da prática pedagógica.
Para tanto, é necessário que o professor planeje e proporcione situações nas quais as interações e \textbf{brincadeiras} sejam eixos articuladores, possibilitando que as \textbf{crianças} se expressem por meio das suas diversas linguagens.
É por meio da interação e do \textbf{brincar} que a infância se caracteriza, que a \textbf{criança} aprende e se desenvolve.
Nessa perspectiva, o documento determina que seis direitos de aprendizagem e desenvolvimento sejam assegurados para as \textbf{crianças} nesta prática, enquanto sujeitos sócio-históricos e culturais.
Sendo assim, tais eixos apontam caminhos que assegurem a compreensão de que as \textbf{crianças} aprendem em situações nas quais possam desempenhar um papel ativo em ambientes que as convidem a vivenciar desafios e a sentirem -se provocadas a resolvê -los, nas quais possam construir significados sobre si, os outros e o mundo social e natural.
Compreendendo garantia de direitos para cumprimento de deveres em busca de uma sociedade alagoana justa, democrática e inclusiva, a BNCC aponta os direitos de aprendizagem e desenvolvimento, que são conviver, \textbf{brincar}, participar, explorar, expressar e conhecer -se.
Ao ingressar no Ensino Fundamental, a \textbf{criança} continua a mesma ; no entanto, a proposta de educação se reconfigura.
As Leis Federais 11.114/2005, 11.274/2006 e 12.796/2013 ( que altera a Lei 9.394/96 em acordo com a Emenda Constitucional 59/2009 ) instituíram uma nova organização do Ensino Fundamental e definiram, em algumas delas, a antecipação do Ensino Fundamental, a ser iniciado aos seis anos de idade, e ampliaram sua duração para nove anos.
Outras ampliaram a obrigatoriedade escolar, estendendo -a dos quatro aos dezessete anos de idade.
Essa organização traz inevitáveis questionamentos e a obrigatoriedade — ao estado conjuntamente aos municípios — na estruturação de orientações que favoreçam a articulação entre as etapas educativas, proporcionando uma transição entre os distintos espaços de socialização da \textbf{criança}, em respeito ao momento das práticas e às concepções de ambas as etapas educacionais, para que venham a ser integradas a partir do reconhecimento de suas diferentes histórias, valores e concepções.
Desse modo, a transição da Educação \textbf{Infantil} para o Ensino Fundamental se constitui em um momento em que se faz necessário o respeito à singularidade das duas etapas, como também um lugar em que a \textbf{criança} precisa transitar alicerçada numa base que garanta seus direitos e vivências de aprendizagem, considerando o passar de uma etapa para outra como parte integrante do processo de desenvolvimento \textbf{infantil} e integral do sujeito.
As DCNEI e a BNCC para a Educação \textbf{Infantil} apontam a necessidade de as instituições de ensino assegurarem essa transição de forma a dar continuidade aos processos de aprendizagem e de desenvolvimento, efetivamente integrando e buscando garantir de fato o que o é de direito à Educação \textbf{Infantil} e à Educação Básica.
Para tanto, esse referencial tem por expectativa apontar elos entre o que se propõe como trabalho de qualidade nas duas etapas.
É imprescindível conhecer o que é próprio de cada etapa e, principalmente, o que é realizado no fazer \textbf{diário} pedagógico, a fim de pensar em estratégias que coloquem a \textbf{criança} em primeiro plano, isto é, ressaltando sua capacidade de aprender e as habilidades a serem conquistadas, considerando o caminho de vivências transcorrido até esse ponto de transição.
Pensar em ações que envolvam ambas as etapas é crucial para que as \textbf{crianças} possam, desde a educação \textbf{infantil}, (re)conhecer os anos iniciais do ensino fundamental como continuidade, espaço de novos desafios e conquistas, onde não se deixa de ser \textbf{criança}, embora o desenvolvimento inerente a cada grupo etário requeira a maturidade que possibilitará a permanência com sucesso dessa \textbf{criança} e seu pleno desenvolvimento.
A transição se caracteriza por mudanças pedagógicas na estrutura educacional, decorrentes principalmente da diferenciação dos componentes curriculares.
Como bem destaca o Parecer CNE/CEB nº 11/2010, “ os estudantes, ao mudarem do professor generalista dos anos iniciais para os professores especialistas dos diferentes componentes curriculares, costumam se ressentir diante das muitas exigências que têm de atender, feitas pelo grande número de docentes dos anos finais ” ( BRASIL, 2010 ).
Realizar as necessárias adaptações e articulações, tanto no 5º quanto no 6º ano, para \textbf{apoiar} os estudantes nesse processo de transição, pode evitar ruptura no processo de aprendizagem, garantindo -lhes maiores condições de sucesso.
Importante destacar que, enquanto sujeitos da ação em idade de \textbf{pré-escola}, a \textbf{criança} precisa ser reconhecida e respeitada a partir de suas especificidades e, para garantir minimamente seu desenvolvimento e formação, é necessário : - Evitar atividades que priorizem a alfabetização e/ou antecipação de conteúdos do EF para \textbf{crianças} em idade de \textbf{pré-escola}, como forma de preparação para a entrada no ensino fundamental ; - Apresentar aos pais a proposta pedagógica da IEI, destacando a importância de atividades \textbf{lúdicas}, com ênfase no \textbf{brincar}, objetivando o desenvolvimento de diversas habilidades e competências para toda a vida ; - Articulação/diálogo entre a Instituição de Educação \textbf{Infantil} ( IEI ) com a Escola de Ensino Fundamental ( EEF ), apresentando proposta/concepção de educação trabalhada de modo que a EEF crie mecanismos de adaptação no primeiro ano ; - Alinhamento da concepção de \textbf{criança}, educação, desenvolvimento e aprendizagem entre as duas etapas da educação básica.

% matched lemmas: adulto, apoiar, bebê, brincadeira, brincar, criança, cuidado, curiosidade, diário, fralda, infantil, lúdico, pequeno, peso, pré-escola
\end{textsample}
